\section{Irrationality by arithmetic specialization}
\label{vana:section}

The modified VAN models admit an arithmetic specialization argument
that is not presently available for the original isotropic-redirection
model.  The key distinction is that their exact characteristic equation
uses a single explicitly controlled transform,
\begin{equation}\label{vana:F}
 F(s)=\int_0^1 x^s e^{x-1}\,dx,\qquad s>-1.
\end{equation}
We first prove the required specialization theorem directly from the
classical Siegel--Shidlovskii theorem.

\begin{theorem}[Transcendence of rational specializations]
\label{vana:rational-specialization}
For every rational number $r>-1$, the number $F(r)$ is transcendental.
\end{theorem}

\begin{proof}
Put $\lambda=r+1>0$ and define
\begin{equation}\label{vana:f}
 f_\lambda(z)=\sum_{n=0}^\infty\frac{z^n}{(\lambda+1)_n}
             =\sum_{n=0}^\infty c_n\frac{z^n}{n!},
 \qquad c_n=\frac{n!}{(\lambda+1)_n},
\end{equation}
where $(u)_0=1$ and $(u)_n=u(u+1)\cdots(u+n-1)$.
The ratio of successive terms proves that $f_\lambda$ is entire.
Substituting $u=1-x$ in~\eqref{vana:F} and integrating the exponential
series gives
\begin{align}
 F(r)
 &=\sum_{n=0}^\infty\frac{(-1)^n}{n!}
      \int_0^1u^n(1-u)^{\lambda-1}\,du\notag\\
 &=\frac1\lambda
      \sum_{n=0}^\infty\frac{(-1)^n}{(\lambda+1)_n}
  =\frac{f_\lambda(-1)}{\lambda}.
 \label{vana:beta-identity}
\end{align}
The interchange is justified by the uniformly absolutely convergent
exponential series and the integrable weight
$(1-u)^{\lambda-1}$.

We verify the arithmetic conditions for $f_\lambda$ to be an
$E$-function.  Write $\lambda=a/b$ with relatively prime positive
integers $a,b$.  Then
\begin{equation}\label{vana:coefficients}
 c_k=\frac{b^k k!}{\prod_{j=1}^k(a+bj)}.
\end{equation}
We claim that the reduced denominators of all $c_0,\ldots,c_n$ divide
the single integer
\begin{equation}\label{vana:denominator}
 D_n=\operatorname{lcm}(1,2,\ldots,a+bn).
\end{equation}
Indeed, if a prime $p$ divides $b$, it divides none of the factors
$a+bj$, so it contributes no denominator.  If $p\nmid b$, the
congruence $a+bj\equiv0\pmod{p^h}$ specifies one residue class
modulo $p^h$.  Hence, for every $k\le n$,
\[
 \#\{1\le j\le k:p^h\mid a+bj\}
 \le \left\lfloor\frac{k}{p^h}\right\rfloor+1.
\]
There are no such factors when $p^h>a+bk$.  Comparing the valuation
of the denominator in~\eqref{vana:coefficients} with Legendre's
formula for $v_p(k!)$ gives
\[
 \max\{0,-v_p(c_k)\}
 \le \left\lfloor\log_p(a+bk)\right\rfloor
 \le \left\lfloor\log_p(a+bn)\right\rfloor
 =v_p(D_n).
\]
This proves the common-denominator assertion.
The elementary bound
$\log\operatorname{lcm}(1,\ldots,M)=O(M)$ therefore yields
$\log D_n=O_\lambda(n+1)$.  Also
\[
 0<c_k=\prod_{j=1}^k\frac{j}{\lambda+j}\le1.
\]
The coefficients are rational, so the same bound holds for all their
algebraic conjugates.  Thus the logarithmic height of
$(c_0,\ldots,c_n)$ is $O_\lambda(n+1)$, as required in the
$E$-function definition.

The coefficient recurrence yields
\begin{equation}\label{vana:first-order}
 z f_\lambda'(z)=(z-\lambda)f_\lambda(z)+\lambda.
\end{equation}
In particular, differentiation gives the homogeneous equation
\[
 z f_\lambda''(z)+(\lambda+1-z)f_\lambda'(z)-f_\lambda(z)=0,
\]
with coefficients in $\mathbb Q[z]$.  This completes the verification
that $f_\lambda$ is an $E$-function.

Every coefficient in~\eqref{vana:f} is nonzero, so $f_\lambda$ is
not a polynomial.  An entire function algebraic over $\mathbb C(z)$
must be a polynomial: a polynomial relation and the usual bound for
roots of a polynomial give polynomial growth outside a disk, after
which Cauchy's estimates force all sufficiently high Taylor
coefficients to vanish.  Consequently $f_\lambda$ is transcendental
over $\mathbb C(z)$.

The two $E$-functions $1,f_\lambda$ satisfy the homogeneous system
\begin{equation}\label{vana:system}
 \frac{d}{dz}
 \begin{pmatrix}1\\f_\lambda(z)\end{pmatrix}
 =
 \begin{pmatrix}
  0&0\\
  \lambda/z&1-\lambda/z
 \end{pmatrix}
 \begin{pmatrix}1\\f_\lambda(z)\end{pmatrix}.
\end{equation}
Its coefficient matrix has common denominator $T(z)=z$.
The Siegel--Shidlovskii theorem in the form of
\cite[Theorem~1.1]{Beukers2006} applies at the nonzero algebraic
point $\xi=-1$, since $\xi T(\xi)=1$.  It gives
\[
 \operatorname{trdeg}_{\overline{\mathbb Q}}
    \overline{\mathbb Q}\bigl(1,f_\lambda(-1)\bigr)
 =
 \operatorname{trdeg}_{\overline{\mathbb Q}(z)}
    \overline{\mathbb Q}(z)\bigl(1,f_\lambda(z)\bigr)
 =1.
\]
Thus $f_\lambda(-1)$ is transcendental.  Equation
\eqref{vana:beta-identity} proves the assertion because
$\lambda$ is a nonzero rational number.
\end{proof}

\begin{theorem}[Irrationality of finite algebraic-parameter VAN roots]
\label{vana:irrationality}
Let $\omega\in\overline{\mathbb Q}\cap[0,\infty)$ be finite.
The unique root $\alpha_\omega\in(0,1)$ of
\begin{equation}\label{vana:root}
 F(\alpha_\omega)
 =\frac{(1-\alpha_\omega)(\alpha_\omega+\omega)}
        {\omega+(1-\omega)\alpha_\omega}
\end{equation}
is irrational.  By Theorem~\ref{van:main}, this is the almost-sure core-growth
exponent of $\operatorname{VAN}(\omega)$.  In
particular, the PAN and DAN exponents are irrational.
\end{theorem}

\begin{proof}
Theorem~\ref{van:main} supplies the existence and uniqueness of
this root in $(0,1)$.  The arithmetic assertion needs
only the root equation.  If $\alpha_\omega$ were rational, its
right-hand side would be algebraic: its denominator satisfies
\[
 \omega+(1-\omega)\alpha_\omega
 =\alpha_\omega+\omega(1-\alpha_\omega)>0.
\]
But Theorem~\ref{vana:rational-specialization} makes its left-hand
side transcendental.  This is a contradiction.  The cases
$\omega=0$ and $\omega=1$ give PAN and DAN, respectively.
\end{proof}

\paragraph{Connection with polynomial exponential integrals.}
For a rational candidate $r=p/q\in(0,1)$ in lowest terms, integration
by parts followed by $x=t^q$ gives the useful normal form
\begin{equation}\label{vana:EIT-form}
 F(p/q)
 =1-\frac pq e^{-1}\int_0^1 x^{p/q-1}e^x\,dx
 =1-p e^{-1}\int_0^1 t^{p-1}e^{t^q}\,dt.
\end{equation}
Thus rational specialization places the problem in the polynomial
exponential-integral setting.  The polynomial amplitude
$p t^{p-1}$ already has degree at most $q-2$.  It is a nonzero
reduced amplitude for the twisted derivative
\[
 \mathcal D_q R=R'+q t^{q-1}R,
 \qquad
 \frac{d}{dt}\bigl(R(t)e^{t^q}\bigr)
       =(\mathcal D_qR)(t)e^{t^q}.
\]
Indeed, for a nonzero polynomial $R$, the highest-degree term of
$\mathcal D_qR$ has degree $\deg R+q-1$, so no nonzero polynomial
of degree at most $q-2$ can be of this form.  Allowing rational $R$
does not create an exception: a finite pole of $R$ gives an
uncancelled pole of one higher order in $R'$, whereas
$q t^{q-1}R$ has no pole of that higher order.  A rational
$R$ with polynomial $\mathcal D_qR$ is therefore polynomial.
This supplies a concrete reduced exponential-integral class to
which an arithmetic specialization theorem can be applied.
For this family, Theorem~\ref{vana:rational-specialization} supplies
that theorem through classical $E$-function theory.

\paragraph{Limits of the conclusion.}
The argument uses rationality of the candidate exponent to obtain an
$E$-function with controlled arithmetic coefficients.  It does not
exclude algebraic irrational roots and therefore does not prove that
$\alpha_\omega$ is transcendental.  The singular model
$\omega=\infty$, whose limiting exponent is $1$, is outside the
theorem.  These are modified local redirection models: no equation
\eqref{vana:root}, and no irrationality or transcendence conclusion,
is asserted for the original isotropic-redirection exponent $\mu$.

\section{A transfer criterion and its scope}
\label{vana:transfer-section}

\begin{proposition}[Transfer with specialization nonvanishing]
\label{vana:transfer}
Let $I\subset(-1,\infty)$ be a real interval.  Suppose an exact
characterization of a target number $\alpha\in I$ gives
\begin{equation}\label{vana:determinant-form}
 D(\alpha)=0,\qquad
 D(s)=\frac{P(s,F(s))}{Q(s,F(s))},
 \qquad P,Q\in\overline{\mathbb Q}[s,X],
\end{equation}
where the denominator is nonzero on $I$.  If
\begin{equation}\label{vana:specialization-nonzero}
 P(r,X)\not\equiv0
 \quad\text{for every }r\in I\cap\mathbb Q,
\end{equation}
then $\alpha$ is irrational.
\end{proposition}

\begin{proof}
If $\alpha=r\in\mathbb Q$, then
\eqref{vana:determinant-form} gives $P(r,F(r))=0$.
Condition~\eqref{vana:specialization-nonzero} makes this a nonzero
polynomial relation over $\overline{\mathbb Q}$ for $F(r)$,
contradicting Theorem~\ref{vana:rational-specialization}.
\end{proof}

The proposition applies, for example, when an exact finite system
has a characteristic determinant of the displayed form.  It is a
criterion to verify, not a claim that the original friend-tree
kernel has such a reduction.  The specialization condition is
essential.  A nonzero bivariate polynomial can vanish identically
as a polynomial in $X$ at a rational parameter: for example,
$P(s,X)=(s-r_0)(X+1)$ vanishes at $s=r_0$ for every value of
$F(r_0)$.  If $P(s,X)=\sum_jP_j(s)X^j$, the needed condition is
precisely that the coefficient polynomials $P_j$ have no common
rational zero in $I$.  This is the nonvanishing step that an exact
elimination must preserve.
